\section{Prospective and retrospective analyses}\label{s:prereg}

The computation mixes proved claims, validated certificates, descriptive
measurements, and statistical models.  We keep those categories separate:
the lower bounds are consequences of the Rabung theorem after machine
verification of its hypotheses, whereas the density discussion of
Section~\ref{s:density} is empirical and model-conditional.

\paragraph{What was fixed prospectively.}
Chunks $1$--$180$ form the training window for the $(3,17)$ comparison.  At
chunk $180$ the repository fixed the baseline count and the three numerical
intensities of Remark~\ref{r:domain}.  A later commit, still before chunk $250$,
fixed the reading schedule $250,300,400,515$.  No parameter was changed after
the cut.  The word \emph{prospective} refers only to the data after chunk
$180$; the discovery of the candidate forms and the choice of their parameters
used the training data and are not called blind.

\paragraph{The original rule.}
The live program compared the cumulative observed increment with the three
means displayed by the frozen functions, identified the closest displayed
mean, and printed Wald-type
discrepancies normalized by $\sqrt{N_{\mathrm{obs}}}$ for two candidates.  This
is the rule that was actually executed at the scheduled reads.  It was not an
exact Poisson test and did not contain the deviance or dispersion analyses now
reported in the paper.  Two manually transcribed standardizations for the
third candidate were erroneous; the counts and forecast functions were not.

\paragraph{Retrospective corrections and sensitivities.}
After completion of the campaign we recomputed the means on the full half-open
chunk intervals, rather than between chunk midpoints, and calculated Pearson
and deviance diagnostics, residual
autocorrelations, a quasi-Poisson scale, an NB2 sensitivity fit, and block
aggregations.  These are labelled post hoc.  They are useful checks on the
specified fixed forecast, but they do not retroactively become part of the
prospective decision rule.  Likewise, freezing $m$ and $\beta$ prevents
holdout refitting but does not make their training-window uncertainty vanish;
we therefore make no family-wide or long-range extrapolative claim from this
comparison.

\paragraph{One trajectory, four scheduled views.}
The four cumulative readings share their early data.  The primary presentation
is therefore the four disjoint increments of Table~\ref{tab:milestones}.  Even
those increments are statistically independent only conditional on the
inhomogeneous-Poisson model.  We retain the cumulative values as an audit trail,
not as four replications, and do not advertise a small conditional tail as a
model-robust significance level.

\paragraph{Scanner revisions.}
The $515$ campaign chunks span eight git revisions and four scanner-source
fingerprints.  The final fingerprint starts at chunk $79$, before the training
cut, so the complete post-freeze comparison uses one scanner source.  This is
the relevant stratification fact for Section~\ref{ss:eta}.  Descriptive
whole-campaign summaries must nevertheless be accompanied by the full revision
table and by tests across the earlier transitions.

\paragraph{Targeted independent checks.}
The historical campaign journal reports that $200$ sampled run-valid $(3,17)$
primes were re-tested through a separate CPU discrete-logarithm walk without a
reported failure.  The present event-list audit does not reproduce that test;
its sampled-prime list and complete raw output must be included and hashed in
the release before the numerical statement is treated as machine-audited.
Even then, the check addresses false survivors, not omissions, prime
enumeration, dispersion, or temporal dependence.

\paragraph{Reproducibility boundary.}
The versioned $1859$-entry list of reported run-valid $(3,17)$ primes is
sufficient to reconstruct all $515$ success counts and the disjoint increments in
Table~\ref{tab:milestones}, conditional on that list.  It does not by itself
re-evaluate the max-run predicate or prove that no run-valid prime was omitted.
The live tracking script additionally reads
\texttt{results/campaign/checkpoint.json}, while the campaign directory is
excluded from ordinary source tracking.  An end-to-end archival claim therefore
requires an immutable copy of the checkpoint, the raw per-chunk outputs, and a
machine-readable revision map in the final data release.  Until those files are
part of the cited release, the git history alone reproduces the narrow
$(3,17)$ event-count analysis, not every campaign-level diagnostic.

This separation is intentional: the prospective record is preserved as it was
executed, later statistical repairs are visible as repairs, and the paper's
conclusions are limited to what the released evidence can support.
